\section{The 40 New City Pairs}

\subsection{Definition and Selection Method}

A ``new city pair'' is defined as an origin-destination pair satisfying three conditions:
\begin{enumerate}
  \item Both the origin city and destination city are within 30 miles of a T1/T1 hub or T1/T2 regional stop in the hub network
  \item The pair currently has no direct intercity bus service with fewer than one transfer (i.e., either no service or service requiring two or more intermediate stops and bus changes)
  \item The pair shows estimated annual demand of at least 5,000 one-way trips, calculated from the gravity model applied to the ACS household data and the FHWA origin-destination commercial traffic matrix \citep{ACS_B08201_2022,FHWA_freight2023}
\end{enumerate}

The second condition --- no current service with fewer than one transfer --- was verified against the January 2026 Greyhound, FlixBus, Megabus, and BoltBus route networks, cross-referenced with the BTS Intermodal Passenger Connectivity Database \citep{BTS_intermodal2023}. Service that operates seasonally or on charter basis is excluded; only scheduled fixed-route service qualifies.

\subsection{The Economic Case for Bus Operator Entry}

A key enabling condition for bus operator entry on new T1 hub city pairs is the managed freight lane PTI of 1.15. At current I-80 PTI of 1.86, transcontinental and long-distance bus scheduling is operationally impossible: a carrier cannot publish a reliable timetable when the corridor PTI implies that 10\% of trips will take 86\% longer than the scheduled time. The PTI 1.15 on managed lanes means that 95\% of trips arrive within 15\% of scheduled time --- a reliability threshold that allows carriers to publish and hold a timetable.

The Gary to Seattle example illustrates this dependency. The T1 express bus on I-90 (Gary, IL to Seattle, WA) is a 28-hour trip at managed lane speeds. At current I-90 PTI of 1.62 (Chicago to Twin Falls segment), the 95th-percentile transit time is 45 hours --- unacceptable for published timetable service. At managed lane PTI 1.15, the 95th-percentile time is 32 hours: long, but schedulable. The managed lane investment makes the city pair viable; the hub infrastructure makes it accessible.

\subsection{Representative New City Pairs}

Table~\ref{tab:pairs} presents 20 representative new city pairs from the 40+ identified in the full analysis. Travel times are estimated at managed lane design speed (62 mph average, including HOS rest stops).

\begin{table}[h]
\centering
\caption{Representative New City Pairs Enabled by T1 Hub Network}
\label{tab:pairs}
\begin{tabular}{llrll}
\toprule
Origin & Destination & Est. Travel Time & Current Service & Nearest Hub(s) \\
\midrule
Atlanta, GA & Jacksonville, FL & 7 hr & 2 transfers via Greyhound & Atlanta; Jacksonville \\
Atlanta, GA & San Antonio, TX & 13 hr & No direct service & Atlanta; San Antonio \\
Atlanta, GA & Richmond, VA & 10 hr & Greyhound (unreliable) & Atlanta; Richmond \\
Jacksonville, FL & Richmond, VA & 11 hr & No direct service & Jacksonville; Richmond \\
Jacksonville, FL & Toledo, OH & 16 hr & No direct service & Jacksonville; Toledo \\
Toledo, OH & Richmond, VA & 9 hr & No direct service & Toledo; Richmond \\
San Antonio, TX & Atlanta, GA & 12 hr & No direct service & San Antonio; Atlanta \\
San Antonio, TX & Jacksonville, FL & 11 hr & No direct service & San Antonio; Jacksonville \\
Sacramento, CA & Gary, IL & 32 hr & 2--3 transfers via Greyhound & Sacramento; Gary \\
Sacramento, CA & Seattle, WA & 14 hr & FlixBus (limited) & Sacramento; Seattle \\
Seattle, WA & Gary, IL & 28 hr & No viable bus option & Seattle; Gary \\
Seattle, WA & Toledo, OH & 30 hr & No viable bus option & Seattle; Toledo \\
Gary, IL & Richmond, VA & 12 hr & No direct service & Gary; Richmond \\
Gary, IL & Jacksonville, FL & 18 hr & No direct service & Gary; Jacksonville \\
Gary, IL & San Antonio, TX & 17 hr & Greyhound (2 transfers) & Gary; San Antonio \\
Boston, MA & Toledo, OH & 13 hr & No direct service & Boston; Toledo \\
Boston, MA & Richmond, VA & 8 hr & Megabus (limited) & Boston; Richmond \\
Richmond, VA & Jacksonville, FL & 11 hr & Greyhound (unreliable) & Richmond; Jacksonville \\
Sacramento, CA & San Antonio, TX & 22 hr & No direct service & Sacramento; San Antonio \\
San Antonio, TX & Toledo, OH & 19 hr & No direct service & San Antonio; Toledo \\
\bottomrule
\end{tabular}
\end{table}

\subsection{The Southeast Cluster: Highest Unmet Demand}

The Southeast cluster --- Atlanta, Jacksonville, Richmond, and the T1/T2 stops along I-20 and I-95 in between --- represents the largest single concentration of unmet intercity bus demand in the hub network. Greyhound's post-2020 restructuring hit the Southeast disproportionately: many small and mid-size cities in Georgia, South Carolina, North Carolina, and Mississippi lost their only intercity bus connection \citep{Amtrak2023}. The T1 hub network, with diamond hubs in Atlanta and Jacksonville and approximately 12 T1/T2 regional stops along I-20 and I-95 in the region, creates the physical substrate for a Southeast intercity bus network that does not currently exist.

The estimated combined annual demand for all Southeast cluster city pairs (Atlanta, Jacksonville, Richmond, plus T1/T2 stops) is 2.3 million one-way trips per year, based on the gravity model calibrated to pre-Greyhound-bankruptcy ridership data and current population estimates. At a 35-seat coach, average 70\% load factor, and 2 round trips per day per route, this demand supports approximately 180 buses per day across the Southeast hub network --- well above the threshold for viable scheduled operations.

\subsection{The Transcontinental Tier: Reliability-Dependent City Pairs}

The transcontinental pairs (Sacramento--Gary, Seattle--Gary, Seattle--Toledo) are categorically different from the Southeast cluster: they are not demand-constrained but reliability-constrained. Bus operators have not offered transcontinental scheduled service since the late 1990s because the PTI on I-80 and I-90 makes timetable commitment impossible. The managed freight lane reduces the PTI to the threshold at which transcontinental scheduling becomes operationally viable. These routes would serve a small but price-sensitive segment of long-distance travelers: students, military families, and low-income travelers for whom air travel is cost-prohibitive \citep{Pisarski2006}. The estimated annual demand on the Sacramento--Gary pair is 28,000 one-way trips --- sufficient for 2 weekly departures at a 50-seat coach, which is the minimum viable service frequency for published timetable operations.
